% ============================================================
%  Prilog: preostale šutnje propisa
%  Jedanaest slučajeva razrađeno je u 6.2, a preostalih sedam ide u
%  prilog kao tablica s citatima.
%
%  S12 je provjeren 2026-08-14: Pravilnik o eRačunu (NN 11/2026, 21 članak)
%    ne sadrži
%    nijednu pojavu riječi "nedostupnost". Dvoznačnost čl. 40. dakle
%    ostaje. Oznaka zatvorena.
%
%  ⚠ S23 je provjeren doslovno na izvoru 2026-08-14 i stranica je
%    arhivirana u rad/izvori/ uz SHA-256, jer nema pročišćen tekst.
% ============================================================
\chapter{Preostale šutnje propisa}
\label{ch:prilog-sutnje}

Odjeljak~\ref{sec:failure-modes} razradio je jedanaest načina otkazivanja koji
imaju posljedicu na zakonsku obvezu. Tijekom izvoda utvrđeno je još sedam mjesta
na kojima propis ne kaže ono što bi izvedba trebala znati. Ona nemaju vlastiti
način otkazivanja, ili ga imaju posredno, pa su ostavljena ovdje uz citate.

Tablica~\ref{tab:sutnje} nije popis primjedbi. Svaka stavka je mjesto na kojem
izvedba mora odlučiti nešto što propis nije odlučio, a takva odluka se ne može
obraniti pozivom na izvor.

\begin{table}[htbp]
\centering
\caption[Preostale šutnje propisa]{Mjesta na kojima propis ne daje podatak koji
izvedba treba. Oznake slijede numeraciju iz izvoda.}
\label{tab:sutnje}
\scriptsize
\setlength{\tabcolsep}{4pt}
\renewcommand{\arraystretch}{1.25}
\begin{tabular}{@{}p{0.7cm}p{5.4cm}p{8.4cm}@{}}
\toprule
& \textsc{Šutnja} & \textsc{Što izvedba mora odlučiti sama} \\
\midrule

S12 & Nedostupnost AMS-a kao servisa nije uređena &
Članak 40.\ pokriva nedostupnost \emph{identifikatora} u AMS-u, dakle slučaj u
kojem se primatelj nije registrirao. Ispad samog servisa gramatički nije
pokriven. Pravilnik tu dvoznačnost ne razrješava. \\

S13 & Nema izlaza za rad bez veze &
Za promet u krajnjoj potrošnji postoje uvezana knjiga računa (čl.~22.) i potvrda
regulatora (čl.~23.). U dijelu o eRačunu ekvivalenta nema. Jedini papirnati
izlaz vezan je isključivo na AMS (čl.~40.~st.~2.). \\

S14 & „Zadovoljavajuća podrška“ nije kvantificirana &
Članak 49.\ stavak 2.\ traži zadovoljavajuću programsku i sklopovsku podršku, a
mjeru ne daje. Izostanak podrške je teži prekršaj
(čl.~72.~st.~1.~t.~10.), pa se sankcionira uvjet koji nije mjerljiv. \\

S19 & Referencijalni integritet reference na prethodni račun &
Nijedno pravilo ne traži da referencirani prethodni račun postoji, da je poslan
ni da je zaprimljen. Schematron traži samo da polje ne bude prazno kada je
skupina prisutna\cite{cen2026validationartefacts}. \\

S20 & Storno nema opisan životni ciklus &
Norma ne poznaje tip dokumenta za poništenje\cite{cen2019en16931}. Hrvatsko
suženje dodaje procesni kod, ali ne određuje semantiku
obrade\cite{pu2026hrcius}. Članak 52.\ stavak 4.\ storniranje ostavlja
mogućnošću, a rok u kojem primatelj smije odbiti nije propisan. \\

S21 & Čuvanje i ponovno izvođenje fiskalizacijskih poruka &
Članak 35.\ propisuje čuvanje eRačuna u izvornom obliku šest godina. O čuvanju
fiskalizacijskih poruka, odgovora i zapisa ne kaže ništa, kao ni o pravu ili
obvezi ponovnog izvođenja tijeka. Europski profil rok čuvanja dokaza
neporecivosti prepušta politici koju nijedan dokument ne
donosi\cite{cef2020as4v115}. \\

S23 & Semantika duplikata živi izvan propisa &
Pravilo da se eRačun pod istim brojem smije poslati samo jednom, obveza
pristupne točke primatelja da vrati potvrdu neovisno o tome je li dokument već
zaprimila i postupanje s tehničkim duplikatom razrađeni su isključivo u rubrici
pitanja i odgovora\cite{pu2026pitanjaodgovori}. Ni Zakon, ni Pravilnik, ni
tehnička specifikacija taj obrazac ne sadrže. \\

\bottomrule
\end{tabular}
\end{table}

Posljednja stavka ima drukčiju narav od ostalih šest i zato je razrađena u
odjeljku~\ref{sec:ogranicenja}. Kod nje propis ne šuti zato što pitanje nije
uočeno. Odgovor postoji i mjerodavan je, ali stoji u dokumentu koji nije propis
i koji se mijenja bez pročišćenog teksta.
